\origpage{350}{N. Robertson, P.D. Seymour / Journal of Combinatorial Theory, Series B 90 (2004) 325--385}
\noindent disc $\Delta \subseteq \hat{\Sigma}$ with either $\Gamma \cap \Delta \subseteq A$ or $\Gamma \cap \Delta \subseteq B$. Let $\Sigma'_1, \Sigma'_2$ be the closures of the two components of $\hat{\Sigma} \setminus U(F')$, where $(\Gamma \cap \Sigma'_1, \Gamma \cap \Sigma'_2) \in \mathscr{T}$. It follows that $(\Gamma \cap \Sigma'_1, \Gamma \cap \Sigma'_2)$ has $\lambda$-order at most that of $(A, B)$.

Suppose that $\Sigma_1 \subseteq \Sigma'_2$. Then $\Sigma'_2$ includes an edge of $A$ (since $\Gamma \cap \Sigma_1 \subseteq A$). Now $(A, B) \in \mathscr{T}$, and $(\Gamma \cap \Sigma'_2, \Gamma \cap \Sigma'_1) \notin \mathscr{T}$, and yet the second separation has order $< ord(\mathscr{T})$, and $\mathscr{T}$ has order $> \frac{1}{2} |E(F)| \ge 1$, and so by the third assertion of Theorem 2.9 of [5], not every edge of $\Gamma \cap \Sigma'_2$ belongs to $A$. Consequently, $\Sigma'_2$ includes an edge of $B$. We deduce that $\Gamma \cap \Sigma'_2 \nsubseteq A$ and $\Gamma \cap \Sigma'_2 \nsubseteq B$. Hence $\Sigma'_2 \ne \Delta$, and so $\Sigma'_1 = \Delta$. Furthermore,
\begin{equation*}
(A \cap \Sigma'_2, B \cup (\Gamma \cap \Sigma'_1)), (A \cup (\Gamma \cap \Sigma'_1), B \cap \Sigma'_2) \in \mathscr{T}
\end{equation*}
because they both have order at most that of $(A, B)$, and $(\Gamma \cap \Sigma'_1, \Gamma \cap \Sigma'_2) \in \mathscr{T}$. Yet if $\Gamma \cap \Sigma'_1 \subseteq A$, the first separation has smaller $\lambda$-order than $(A, B)$, while if $\Gamma \cap \Sigma'_1 \subseteq B$ the second does (for in both cases the first term of the $\lambda$-order does not increase, and the second term strictly decreases). In either case this is contrary to our choice of $(A, B)$; and one of these occurs since $\Sigma'_1 = \Delta$, a contradiction. Thus, $\Sigma_1 \nsubseteq \Sigma'_2$.

Now no edge of $F'$ lies in $\Sigma_1 \setminus U(F)$, because every edge of $F'$ is split; and so $\Sigma_1 \subseteq \Sigma'_1$ and $\Sigma'_2 \subseteq \Sigma_2$. Hence $(\Gamma \cap \Sigma'_1, \Gamma \cap \Sigma'_2)$ has $\lambda$-order at least that of $(A, B)$, from the choice of $(A, B)$. Since we already shown the reverse inequality, it follows that the $\lambda$-orders are equal and so from 8.2, $(A, B) = (\Gamma \cap \Sigma'_1, \Gamma \cap \Sigma'_2)$. Hence $F'$ is $\lambda$-linked to $\mathscr{T}$, and (from the choice of $(A, B)$) $|E(F')| \le |E(F)|$, and if equality holds then $\lambda(F') \le \lambda(F)$. This proves 8.7. $\square$

An \textit{arm} $A$ of a painting $\Gamma$ is a pair $(A^-, \pi(A))$, where $A^-$ is a subhypergraph of $\Gamma$ and $\pi(A)$ is a march in $V(A^-)$, with the property that there exists $B \subseteq \Gamma$ such that $(A^-, B)$ is a separation of $\Gamma$ and $V(A^- \cap B) = \bar{\pi}(A)$. (In other words, for each edge $e$ of $\Gamma$ not in $E(A^-)$, $\bar{\pi}(A)$ contains every end of $e$ in $V(A^-)$.) In this case, we call $B$ the \textit{complement} of $A$ (it is unique) and write $B = A^c$. We define $V(A) = V(A^-), E(A) = E(A^-)$. We say that $A$ is $\lambda$-\textit{linked to} $\mathscr{T}$ if $(A^-, A^c)$ is $\lambda$-linked to $\mathscr{T}$. The \textit{order} of $A$ is $|\bar{\pi}(A)|$, that is, the order of $(A^-, A^c)$.

A \textit{rooted location} in $\Gamma$ is a set $\mathscr{L}$ of arms, such that that if $A_1, A_2 \in \mathscr{L}$ with $A^-_1 \ne A^-_2$ then $A^-_1 \subseteq A^c_2$. Its \textit{order} $ord(\mathscr{L})$ is the maximum order of its members (or zero, if $\mathscr{L} = \emptyset$). We define $\mathscr{L}^- = \{(A^-, A^c) : A \in \mathscr{L}\}$. A rooted location $\mathscr{L}$ with $\mathscr{L}^- \subseteq \mathscr{T}$ is $\lambda$-\textit{linked to} $\mathscr{T}$ if each of its members is $\lambda$-linked to $\mathscr{T}$. A rooted location $\mathscr{L}$ \textit{$\theta$-isolates} a tangle $\mathscr{T}$ if $\mathscr{L}^- \subseteq \mathscr{T}, ord(\mathscr{L}) < \theta \le ord(\mathscr{T})$, and for every $A \in \mathscr{L}$ and for every tangle $\mathscr{T}'$ of order $\ge \theta$ with $(A^c, A^-) \in \mathscr{T}'$, the $(\mathscr{T}, \mathscr{T}')$-distinction $(C, D)$ satisfies $C \subseteq A^-$ and $A^c \subseteq D$. It follows by applying theorem 7.1 of [10] to $\mathscr{L}^-$ that

\textbf{8.8.} \textit{If $\mathscr{T}$ is a tangle in $\Gamma$, and $\mathscr{L}$ is a rooted location with $\mathscr{L}^- \subseteq \mathscr{T}$, and $ord(\mathscr{L}) < \theta \le ord(\mathscr{T})$, and $\mathscr{L}$ is $\lambda$-linked to $\mathscr{T}$, then $\mathscr{L}$ $\theta$-isolates $\mathscr{T}$.}
